\chapter{Contexto general del repositorio}

El código se organiza en tres bloques funcionales:

\begin{itemize}[leftmargin=1.5cm]
  \item \texttt{src/algoritmos}: implementaciones de algoritmos de grafos reutilizables (coloreo, caminos mínimos, flujo, MST, recorridos).
  \item \texttt{src/estructuras}: tipos y estructuras de soporte (\texttt{Vertex}, \texttt{Edge}, \texttt{Graph}, proyecciones, \texttt{PriorityQueue}).
  \item \texttt{src/universidad-proyecto}: solución aplicada con cargadores de datos, constructor de grafo de conflictos, solucionadores, API y frontend.
\end{itemize}

La separación permite reutilizar algoritmos y estructuras en otros problemas y facilita la comprobación experimental de distintas heurísticas sobre la misma instancia.

